% !TEX root = ../main.tex
% Diagramma UML di sequenza della pipeline statica run_pipeline()
% dell'orchestratore: scansioni, correlazione e punteggio. Il secondo
% argomento di correlate (segnalazioni a runtime) e' vuoto.
% NB: \mess usa il testo del messaggio come nome di nodo TikZ: niente
% parentesi, virgole o macro nel terzo argomento.
\linespread{1.0}\selectfont
% Ridotto alla larghezza del testo solo se la eccede.
\resizebox{\ifdim\width>\linewidth\linewidth\else\width\fi}{!}{%
\begin{sequencediagram}
  % passo verticale fra i livelli (default 0.6cm)
  \renewcommand{\unitfactor}{0.6}
  \tikzstyle{inststyle}+=[font=\footnotesize, rounded corners=0pt, fill=white]
  \tikzset{every node/.style={font=\footnotesize}}
  \newthread[white]{orc}{Orchestratore}
  \newinst[1.2]{sc}{CheckovScannerAdapter}
  \newinst[1.6]{eng}{RiskCorrelationEngine}
%
  \begin{sdblock}{loop}{per ogni scanner registrato}
    \begin{call}{orc}{scan(target\_dir)}{sc}{list[Finding]}
    \end{call}
  \end{sdblock}
  \postlevel
%
  \begin{call}{orc}{correlate(static, [])}{eng}{list[Finding]}
    \begin{callself}[2]{eng}{\shortstack[l]{indicizza per chiave; per ogni risorsa sceglie\\ la rappresentante (natura, poi severità)\\ e aggrega gli altri controlli}}{}
    \end{callself}
  \end{call}
%
  \begin{sdblock}{loop}{per ogni finding correlato}
    \begin{call}{orc}{calculate\_risk\_score(finding)}{eng}{score}
    \end{call}
  \end{sdblock}
  \postlevel
\end{sequencediagram}%
}
